\documentclass[a4paper,landscape,9pt,twocolumn]{article}

% ---- Language / Font ----
\usepackage[utf8]{inputenc}
\usepackage[T1]{fontenc}
\usepackage{lmodern}
\usepackage{kotex}
\usepackage{fontspec}
\usepackage{mathtools}
\usepackage{graphicx}

\setmainhangulfont{Noto Serif CJK KR}

% ---- Page layout ----
\usepackage{geometry}
\geometry{
  left=14mm, right=14mm,
  top=14mm, bottom=14mm,
  headsep=3mm
}

% ---- Header / Footer ----
\usepackage{fancyhdr}
\usepackage{lastpage}

\pagestyle{fancy}
\fancyhf{}
\lhead{\small\textbf{Semi-Game Cup 4}}
\rhead{\small\texttt{kdy40929}}
\cfoot{\small \thepage\ / \pageref{LastPage}}
\renewcommand{\headrulewidth}{0.4pt}
\renewcommand{\footrulewidth}{0pt}

% ---- Section style ----
\usepackage{titlesec}
\titleformat{\section}{\large\bfseries}{\thesection.}{0.5em}{}
\titleformat{\subsection}{\normalsize\bfseries}{\thesubsection.}{0.5em}{}
\titleformat{\subsubsection}{\small\bfseries}{\thesubsubsection}{0.5em}{}

% 제목 간격(위/아래): left, before, after
\titlespacing*{\section}{0pt}{0.6em}{0.2em}
\titlespacing*{\subsection}{0pt}{0.4em}{0.2em}
\titlespacing*{\subsubsection}{0pt}{0.3em}{0.1em}

% ---- Compact spacing ----
\usepackage{enumitem}
\setlist{nosep}
\setlength{\parskip}{0.1em}
\setlength{\parindent}{0pt}

% 두 단에서 간격 조금 줄이기 (원하면 더 조절)
\setlength{\columnsep}{6mm}

% ---- Code blocks ----
\usepackage[highlightmode=immediate]{minted}
\usepackage{xcolor}

\setminted{
  fontsize=\fontsize{8}{9.5}\selectfont,
  breaklines,
  breakanywhere,
  tabsize=2,
  autogobble
}

\usemintedstyle{manni}

\begin{document}
\raggedbottom

% =========================
% Table of Contents
% =========================
\section*{kdy40929's Reference List v1.0 - 2026.01.12}
\vspace{-2mm}
\renewcommand{\contentsname}{}
\tableofcontents

\newpage

% =========================
% Main Contents
% =========================

\section{Math}
\subsection {Linear Algebra}
\subsubsection{Matrix Operations}
\inputminted{python}{codes/math/matrix.py}
\subsubsection{Gauss Elimination}
\inputminted{python}{codes/math/gauss.py}
\subsubsection{XOR basis}
\inputminted{python}{codes/math/xor_basis.py}

\subsection{Number Theory}
\subsubsection{Euler Phi}
\inputminted{python}{codes/math/euler_phi.py}
\subsubsection{Miller labin \& Pollard Rho}
\inputminted{python}{codes/math/miller_labin.py}
\subsubsection{Combination with Lucas Theorem}
\inputminted{python}{codes/math/lucas.py}

\subsection{Geometry}
\subsubsection{Geometry}
\inputminted{python}{codes/math/geo.py}

\subsection{Others}
\subsubsection{FFT}
\inputminted{python}{codes/math/fft.py}
\subsubsection{Floor Sum}
\inputminted{python}{codes/math/floorsum.py}

\section{Data Structures}
\subsection{Deque}
\subsubsection{Deque Trick}
\inputminted{python}{codes/ds/deque.py}
\subsection{Segment Tree}
\subsubsection{Fenwick Tree}
\inputminted{python}{codes/ds/fenwick.py}
\subsubsection{Non-recurisve Range Update Range Query}
\inputminted{python}{codes/ds/lazyseg.py}
\subsubsection{Merge Sort Tree}
\inputminted{python}{codes/ds/mst.py}
\subsubsection{2D Segment Tree}
\inputminted{python}{codes/ds/2dseg.py}
\subsubsection{Maximum Subarray Sum Segment Tree}
\inputminted{python}{codes/ds/mineseg.py}
\subsection{Trie}
\subsubsection{1D Trie}
\inputminted{python}{codes/ds/trie.py}

\section{Graphs}
\subsection{Shortest Path}
\subsubsection{Dijkstra}
\inputminted{python}{codes/graphs/dijkstra.py}
\subsubsection{Bellman Ford}
\inputminted{python}{codes/graphs/bellman_ford.py}
\subsubsection{SPFA}
\inputminted{python}{codes/graphs/spfa.py}
\subsubsection{Floyd Warshall}
\inputminted{python}{codes/graphs/floyd.py}
\subsection{SCC and similar}
\subsubsection{SCC}
\inputminted{python}{codes/graphs/scc.py}
\subsubsection{2-SAT}
\inputminted{python}{codes/graphs/2sat.py}
\subsubsection{Articulation}
\inputminted{python}{codes/graphs/articulation.py}
\subsection{Matching \& Flow}
\subsubsection{Bipartite Matching - Hopcroft Karp}
\inputminted{python}{codes/graphs/hopcroft.py}
\subsubsection{Dinic Algorithm}
\inputminted{python}{codes/graphs/dinic.py}
\subsubsection{Dinic with low memory}
\inputminted{python}{codes/graphs/dinic_opt.py}
\subsubsection{MCMF}
\inputminted{python}{codes/graphs/mcmf.py}
\subsection{Trees}
\subsubsection{LCA}
\inputminted{python}{codes/graphs/lca.py}
\subsubsection{HLD}
\inputminted{python}{codes/graphs/hld.py}

\section{DP}
\subsection{Basic DP}
\subsubsection{LIS}
\inputminted{python}{codes/dp/lis.py}
\subsection{DP optimization}
\subsubsection{Convex Hull Trick}
$dp[i] = \min_{j < i} (dp[j] + A[j]*B[i])$ 형태의 관계식을 $A[j] = a, B[i] = x, dp[j] = b$로 하여 $y = ax+b$ 꼴의 직선들의 볼록 껍질을 만들어 해결한다.
$A[j]$의 단조성이 보장되어야 한다.\\
\inputminted{python}{codes/dp/cht.py}

\section{String}
\subsection{Pattern Matching}
\subsubsection{KMP}
\inputminted{python}{codes/string/kmp.py}
\subsubsection{Aho Corasick}
\inputminted{python}{codes/string/corasick.py}
\subsection{Substring Algorithm}
\subsubsection{Rolling Hash}
\inputminted{python}{codes/string/hashing.py}
\subsubsection{Z}
\inputminted{python}{codes/string/z.py}
\subsubsection{Manacher}
\inputminted{python}{codes/string/manacher.py}
\subsubsection{Suffix array \& LCP array}
\inputminted{python}{codes/string/salcp.py}

\section{Misc}
\subsection{Useful Codes}
\subsubsection{Fast I/O}
\inputminted{python}{codes/misc/fastio.py}
\subsubsection{Coordinate Compression}
\inputminted{python}{codes/misc/compress.py}
\subsection{Others}
\subsubsection{Graph Theorem}
쾨닉의 정리: 이분 그래프에서 (최대 매칭) = (최소 정점 커버) \\
이분 그래프에서 (정점 수) = (최대 매칭) + (최대 독립집합 크기) \\
홀의 결혼 정리: 이분 그래프에서 좌측 집합 $L$의 모든 부분집합 $S \subset L$에 대해 $|N(S)| \ge |S|$이면 $L$에 대한 완전 매칭 존재
\subsubsection{Useful Number}
NTT Friendly Prime \\
$ 998244353 = 119 \times 2^{23} + 1 $, Primitive Root = $3$ \\
$ 985661441 = 235 \times 2^{22} + 1 $, Primitive Root = $3$ \\
Useful Prime Number \\
$10^{8} + 7$, $10^{9} + 7$, $10^{9} + 9$, $10^{11} + 3$, $10^{12} + 39$ $10^{18} + 3$ \\

\subsubsection{Useful Sequences}
Fibonacci $F_n$ \\
$F_0 = 0$, $F_1 = 1$, $F_n = F_{n-1} + F_{n-2}$ \\
$0$, $1$, $1$, $2$, $3$, $5$, $8$, $13$, $21$, $34$, $55$, $89$, $144$, $233$, $377$, $610$, $987$, $1597$, $2584$ $\cdots$
\newline

Lucas $L_n$ \\
$L_0 = 2$, $L_1 = 1$, $L_n = L_{n-1} + L_{n-2}$ \\
$2$, $1$, $3$, $4$, $7$, $11$, $18$, $29$, $47$, $76$, $123$, $199$, $322$, $521$, $843$, $1364$, $2207$ $\cdots$
\newline

Tribonacci $T_n$ \\
$T_0 = 0$, $T_1 = 0$, $T_2 = 1$, $T_n = T_{n-1} + T_{n-2} + T_{n-3}$ \\
$0$, $0$, $1$, $1$, $2$, $4$, $7$, $13$, $24$, $44$, $81$, $149$, $274$, $504$, $927$, $1705$ $\cdots$
\newline

Catalan $C_n$ \\
$C_0 = 1$, $C_n = \frac{1}{n+1}\binom{2n}{n}$ \\
$1$, $1$, $2$, $5$, $14$, $42$, $132$, $429$, $1430$, $4862$, $16796$, $58786$, $208012$, $742900$ $\cdots$
\newline

Central Binomial Coefficient $B_n$ \\
$B_n = \binom{2n}{n}$ \\
$1$, $2$, $6$, $20$, $70$, $252$, $924$, $3432$, $12870$, $48620$, $184756$, $705432$, $2704156$ $\cdots$
\newline

Partition Number $p(n)$ \\
$p(0)=1$, $p(n)=\sum_{k\neq 0} (-1)^{k-1} p!\left(n-\dfrac{k(3k-1)}{2}\right)$ \\
$1$, $1$, $2$, $3$, $5$, $7$, $11$, $15$, $22$, $30$, $42$, $56$, $77$, $101$, $135$, $176$, $231$ $\cdots$
\newline

Bell $B_n$ \\
$1$부터 $n$까지 자연수로 이루어진 집합을 분할하는 방법의 수 \\
$B_0 = 1$, $B_{n+1} = \sum_{k=0}^{n} \binom{n}{k} B_k$ \\
$1$, $1$, $2$, $5$, $15$, $52$, $203$, $877$, $4140$, $21147$, $115975$, $678570$ $\cdots$
\newline

Partition number $p(n,k)$ \\
정수 $n$을 정확히 $k$개의 양의 정수의 합으로 나타내는 방법의 수 (순서는 무시) \\
$p(0,0)=1$, $p(n,0)=0\ (n>0)$, $p(0,k)=0\ (k>0)$ \\
$p(n,k)=p(n-1,k-1)+p(n-k,k)$ \\
\newline

Stirling (1st kind) $s(n, k)$ \\
$x^{\underline{n}}$, $x$의 $n$ 하강 계승의 $x^{k}$ 계수  \\
$s(0,0)=1$, $s(n,0)=0\ (n>0)$, $s(n,k)=s(n-1,k-1)-(n-1)s(n-1,k)$ \\
\newline

Stirling (2nd kind) $S(n,k)$ \\
$1$부터 $n$까지의 자연수로 이루어진 집합을 $k$개로 분할하는 방법의 수 \\
$S(0,0)=1$, $S(n,0)=0\ (n>0)$, $S(n,k)=S(n-1,k-1)+kS(n-1,k)$ \\
\newline



\end{document}
